% questions.tex —— 1977 年·江西卷 · 题目内容片段（无答案）
% 由 tools/generate_projects.py 依来源总账自 1977/1977_数学真题_江西_试题.md 生成，勿手改（改源 MD 后重跑）


\gkpart{数学试卷（江西卷）}

\begin{gktip}
\textbf{来源说明}：本卷依据固定源 PDF《1977江西.pdf》（11 页扫描）逐题转录。固定源并非空白试卷，而是一册题解排印本（书页 40—49，页脚印有「题解—40」等页码），题面、解答与〔评分标准〕连排，卷首标题作「江西省高等学校招生数学试题」。正式题为第一—六题，末尾〔附加题〕第七题（20 分）不计入正式题量；全卷无文理分卷说明。源卷第一大题第（7）—（15）小项括号内印有答案、第（1）—（6）小项括号留空，本卷题面一律留空，括号内答案见答案要点文件。原卷无插图（第六题解答注明「图自画」）。
\end{gktip}

\par\medskip

\begin{question}
\gkti{1}\textbf{（20 分）} 下列各题将答案填在括号内：

\begin{gktip}
\textbf{题号说明}：源卷第一题的 15 个小项印作「1.」—「15.」，本库规范要求小问用全角括号编号，故重排为（1）—（15），未增删题干内容。
\end{gktip}

（1） $f(x)=3x^3-5x^2+7x-5$，则 $f(0)=$\mbox{（　　）}；

（2） $\operatorname{arctg}1=$\mbox{（　　）}；

（3） $\log_a a^2=$\mbox{（　　）}；

（4） 球的面积是大圆面积的\mbox{（　　）}倍；

（5） 若 $a<0$，则 $\sqrt{a^2}\div a=$\mbox{（　　）}；

（6） $\dfrac12+\dfrac{\sqrt3}{2}i$ 的模是\mbox{（　　）}，幅角是\mbox{（　　）}，三角函数式是\mbox{（　　）}；

（7） 若 $\alpha$ 为锐角，又 $\sin 2\alpha=1$，则 $\sin\alpha=$\mbox{（　　）}；

（8） 若 $a$ 大于 $b$、$c$ 小于零，则 $ac$\mbox{（　　）}于 $bc$；

（9） 三角形的重心是三条\mbox{（　　）}线的交点；

（10） 若 $C_n^3=C_n^5$，则 $n=$\mbox{（　　）}；

（11） 直线的方程为 $y=-3x+2$，则它的斜率是\mbox{（　　）}；

（12） 正圆锥底半径长 $3$ 公分，母线长 $5$ 公分，则高是\mbox{（　　）}公分；

（13） $y=2\sin\left(3x+\dfrac{\pi}{6}\right)$ 的振幅是\mbox{（　　）}，周期是\mbox{（　　）}；

（14） $\dfrac{\sqrt3+\sqrt2}{\sqrt3-\sqrt2}$ 与 $\dfrac{\sqrt3-\sqrt2}{\sqrt3+\sqrt2}$ 的等差中项是\mbox{（　　）}，等比中项是\mbox{（　　）}；

（15） 圆的半径为 $R$，内接正方形的边长为\mbox{（　　）}，内接正三角形边长为\mbox{（　　）}。

\begin{gktip}
\textbf{源注}：原卷「圆」字排作「园」（如「正园锥」「园的半径」），本文统一转写为「圆」。原卷第（7）—（15）小项括号内排有答案，本卷一律留空。
\end{gktip}

\end{question}
\begin{question}
\gkti{2}\textbf{（15 分）} 若方程 $5x^2-10x\cos\alpha+7\cos\alpha+6=0$ 的两根相等，试求两邻边之和为 $6$，且夹角为 $\alpha$ 的平行四边形的最大面积。

\end{question}
\begin{question}
\gkti{3}\textbf{（15 分）} 若

\nopagebreak[4]
\[
5\left|\frac{a+2b-3}{b+1}\right|+3(a+3b)^2=0,
\]

求无穷数列 $ab$，$b$，$\dfrac{b}{a}$，$\cdots\cdots$ 的和。

\begin{gktip}
\textbf{源注}：源卷该数列第三项 $\dfrac{b}{a}$ 前印有一短横状墨迹，形似负号；按源解答所列数列前三项 $-27$、$-3$、$-\dfrac13$（公比 $\dfrac19$）判读应为 $\dfrac{b}{a}$，本文按此转录并留注。
\end{gktip}

\end{question}
\begin{question}
\gkti{4}\textbf{（20 分）} 用 $0$、$2$、$5$、$7$、$9$ 诸数字组成没有重复数字的四位数中

（1） 大于 $5000$ 的有多少个？

（2） 是 $4$ 的倍数的有多少个？

\end{question}
\begin{question}
\gkti{5}\textbf{（15 分）} 三角形的三边为 $a$，$b$，$c$，它们的对角分别为 $A$、$B$、$C$，若 $a^2=b(b+c)$，求证 $A=2B$。

\begin{gktip}
\textbf{源注}：原卷此行铅字排印作「b(b+C)」「4=2B」——系该排印本字形与大小写混排所致（斜体 A 形近「4」、边 c 排作大写 C），按题干「对角分别为 A、B、C」的对应关系分别转写为 b(b+c) 与 A=2B。
\end{gktip}

\end{question}
\begin{question}
\gkti{6}\textbf{（15 分）} 设等腰 $\triangle ABC$ 的底边 $BC$ 在平面 $P$ 上，它的底角为 $75^\circ$，顶点 $A$ 到平面 $P$ 的距离 $AD=4$ 公分，$\triangle ABC$ 所在平面与平面 $P$ 成 $60^\circ$ 的角，求棱锥 $A-BCD$ 的体积。

\begin{gktip}
\Needspace{0.25\textheight}
\textbf{附加题}
\end{gktip}

\textbf{附加题 7．（20 分）} 已知抛物线 $y^2=x$ 上两点 $A(1,1)$ 和 $B(4,-2)$，求过 $A$、$B$ 的直线 $AB$ 与抛物线 $y^2=x$ 所围成图形的面积。

\end{question}